\section{Wstęp}
\label{sec:wstep}

\subsection*{Cel projektu}

Celem projektu jest opracowanie narzędzia wspomagającego planowanie trasy
turystycznej na zadanej mapie. Turysta wyrusza z~ustalonego punktu \emph{start},
porusza się po siatce komórek o~zróżnicowanych kosztach przejścia i~ma za zadanie
zakończyć podróż w~punkcie \emph{end}, odwiedzając po drodze atrakcje w~taki sposób,
aby zmaksymalizować łączną wartość zdobytych atrakcji przy jednoczesnym spełnieniu
ograniczeń:

\begin{itemize}
    \item budżetowego (łączny koszt biletów wstępu do odwiedzonych atrakcji nie może
          przekroczyć założonego budżetu),
    \item czasowego (łączny czas: ruch po siatce plus czas zwiedzania nie może
          przekroczyć dostępnego budżetu czasu),
    \item dotyczącego liczności typów atrakcji (dla każdego typu zdefiniowane są
          minimalna i~maksymalna liczba atrakcji, które trzeba lub można odwiedzić).
\end{itemize}

Dodatkowo, ze względów modelowych, przyjęto twarde wymaganie, aby trasa
\emph{nie używała dwukrotnie tej samej krawędzi} (przejścia między dwoma sąsiednimi
polami) --- co odpowiada intuicji, że turysta nie chce przechodzić tą samą wąską
ścieżką w~obie strony.

\subsection*{Założenia projektu}

\begin{enumerate}
    \item Mapa jest dyskretną siatką $H \times W$ komórek. W~każdej komórce
          $(r,c)$ zdefiniowana jest waga $w_{r,c} \in \N_{+}$ reprezentująca koszt
          ,,bycia'' w~tej komórce (im wyższa waga, tym trudniejszy/dłuższy teren).
    \item Ruch w~siatce odbywa się w~jednym z~ośmiu kierunków
          (4 osiowe + 4 diagonalne). Koszt przejścia z~komórki $(r,c)$ do sąsiada
          wynosi $w_{r,c}$ pomnożone przez $1{,}0$ dla ruchu osiowego oraz $1{,}4$
          dla ruchu diagonalnego (przybliżenie $\sqrt{2}$).
    \item Atrakcje to wyróżnione komórki o~zadanej wartości, koszcie biletu, czasie
          zwiedzania i~typie. Wizyta polega na wejściu na komórkę atrakcji; ponowne
          odwiedzenie tej samej atrakcji nie daje żadnych dodatkowych punktów.
    \item Algorytm operuje w~paradygmacie \emph{generowania populacji rozwiązań
          dopuszczalnych} --- każde rozwiązanie w~populacji końcowej musi spełniać
          wszystkie ograniczenia (budżet, czas, liczności typów, brak powtórzonych
          krawędzi).
    \item Komponentem dostarczającym rozwiązanie początkowe i~jednocześnie
          mechanizmem poszukiwania jest meta-heurystyka \emph{Artificial Bee Colony}
          (ABC).
    \item Aplikacja składa się z~modułów: \texttt{map\_generator} (generator instancji),
          \texttt{parser}/\texttt{load\_json} (obsługa formatu danych),
          \texttt{abc\_algorithm} (rdzeń algorytmu), \texttt{initial\_population}
          (skrypt uruchomieniowy generujący populację) oraz \texttt{visualize\_paths}
          (wizualizacja). Dodatkowo dostępna jest aplikacja webowa
          (\texttt{app.py}, Streamlit) służąca jako interaktywny edytor mapy.
\end{enumerate}

\subsection*{Zakres dokumentu}

Dokument prezentuje formalny opis problemu i~jego model matematyczny
(rozdział~\ref{sec:zagadnienie}), szczegółowo omawia zaprojektowaną metodę
rozwiązania (rozdział~\ref{sec:algorytmy}), opis aplikacji wraz z~instrukcją
uruchomienia (rozdział~\ref{sec:aplikacja}), zarys planowanych eksperymentów
(rozdział~\ref{sec:eksperymenty}), wnioski i~kierunki rozwoju
(rozdział~\ref{sec:podsumowanie}) oraz literaturę i~dodatek z~podziałem prac.
